\section{Catene di Markov e PageRank}

\subsection{Catena di Markov}

Stati $S_j$ (prima della transizione), $S_i^\to$ (dopo la transizione). Matrice di transizione $T$ (omogenea: indipendente dal tempo)
$$ \boxed{ t_{ij} = P(S_i^\to \mid S_j) } \qquad \sum_i t_{ij} = 1 \quad \forall j $$

Distribuzione $p$ con $p_i = P(S_i)$ (probabilità totale):
$$ \boxed{ p' = T \, p } \qquad \boxed{ p(t+n) = T^n \, p(t) } $$


\subsection{Distribuzione stazionaria}

\textbf{Perché un eigenvektor?} 

Dopo $n$ passi $p(n) = T^n p_0$. Se per $n \to \infty$ la distribuzione converge a $p$, un passo in più non cambia nulla: la probabilità ``dopo'' è uguale a quella ``prima'', $P(S_i^\to) = P(S_i)$, cioè $p' = p$. \\
Con $p' = T\,p$ segue $T\,p = 1 \cdot p$ \textrightarrow\ \textbf{eigenvektor con eingwert $1$}


\begin{tabular}{ll}
	stazionaria			& $T \, p = p$ \textrightarrow\ eigenvektor con eingwert $1$ \\
	distribuzione limite	& $\displaystyle \lim_{n \to \infty} T^n p_0 = p$ (è stazionaria)
\end{tabular}



\subsection{Google PageRank}

$S_j$ = visitatori sulla pagina $j$, $S_i^\to$ = visitatori sulla pagina $i$ dopo la navigazione

\textbf{Matrice dei link} $H$ \quad ($h_{ij} = P(S_i^\to \mid S_j)$ = probabilità di andare dalla pagina $j$ alla pagina $i$, somma delle colonne $1$):
$$ \boxed{ H = \begin{pmatrix}
	P(S_1^\to \mid S_1) & \cdots & P(S_1^\to \mid S_N) \\
	\vdots & \ddots & \vdots \\
	P(S_N^\to \mid S_1) & \cdots & P(S_N^\to \mid S_N)
\end{pmatrix} } $$

Libero arbitrio $F$ (l'utente non segue un link), peso $\omega$ per la navigazione tramite link:
$$ \boxed{ G = \omega \, H + (1-\omega) \, \frac{1}{N} A } \qquad A = \begin{pmatrix} 1 & \cdots & 1 \\ \vdots & & \vdots \\ 1 & \cdots & 1 \end{pmatrix} $$

\textbf{PageRank} = autovettore $p$ di $G$ con autovalore $1$:
$$ \boxed{ G \, p = p } \qquad p = \lim_{n \to \infty} G^n p_0 $$